% Part: proof-theory
% Chapter: natural-deduction
% Section: translation-G2i

\documentclass[../../../include/open-logic-section]{subfiles}

\begin{document}

\olfileid{pt}{nat}{tgi}

\olsection{\Log{G2i} থেকে \Log{N2i}-তে অনুবাদ}

\Log{G2i}-এর সিকোয়েন্টের বাঁ অংশ হলো
!!{formula}s-এর বহুসেট~$\Gamma$; কিন্তু
\Log{N2i}-এর সিকোয়েন্টের বাঁ অংশ হলো
চিহ্নযুক্ত সূত্রের সেট~$\Gamma'$, যেখানে
প্রতিটি চিহ্ন মাত্র একবার ঘটে। চিহ্নযুক্ত
!!{formula}s-এর সেট~$\Gamma'$-কে বহুসেট~$\Gamma$-র
\emph{সঙ্গত} বলব যদি প্রত্যেক !!{formula}~$!A$-এর
জন্য $\Gamma'$-তে $x : !A$ রূপের !!{element}s-এর
সংখ্যা $\Gamma$-তে~$!A$-এর পুনরাবৃত্তি-সংখ্যার
(অর্থাৎ~$!A$-এর অনুলিপির সংখ্যার) বেশি না হয়।

\begin{prop}\ollabel{prop:G2i-to-N2i}
যদি $\Log{G2i} + \Cut \Proves \Gamma \Sequent \Delta$,
তাহলে $\Log{N2i} \Proves \Gamma' \Sequent !C$,
যেখানে $\Delta$ ফাঁকা হলে $!C \ident \lfalse$,
অন্যথায় $\Delta = \{!C\}$; এবং $\Gamma'$
হলো~$\Gamma$-র সঙ্গত।
% BN-SRC-876 TARGET-CORRECTION-BEGIN
\par\smallskip\noindent\textnormal{\small\textbf{পাঠ-সংশোধন (BN-SRC-876):} অফাঁকা ডান অংশের একমাত্র সূত্রই সিদ্ধান্তের সূত্র।
মূলপাঠে অন্য একটি নাম থাকায় সিদ্ধান্তের সূত্রটি অনির্ধারিত ছিল;
দুই স্থানে একই নাম রাখা হয়েছে।}
% BN-SRC-876 TARGET-CORRECTION-END
\end{prop}

\begin{proof}
  দেখাই, $\pi$ যদি \Log{G2i}+\Cut-এ
  $\Gamma \Sequent \Delta$-র !!a{proof} হয়,
  তাহলে $\Gamma$-র সঙ্গত কোনো~$\Gamma'$ এবং
  \Log{N2i}-তে $\Gamma' \Sequent !C$-এর
  !!a{proof}~$\delta$ আছে। $n = \pheight{\pi}$-র
  উপর আরোহ করি।
  % BN-SRC-559 TARGET-CORRECTION-BEGIN
  \par\smallskip\noindent\textnormal{\small\textbf{উৎস-সংশোধন (BN-SRC-559):} প্রমাণের শুরুতে উৎসে অনির্দিষ্ট G2ci লেখা; প্রস্তাবনা ও সারা অধ্যায় অনুযায়ী G2i নেওয়া হয়েছে।}
  % BN-SRC-559 TARGET-CORRECTION-END
% BN-SRC-877: redo restores G2i+Cut scope of the proposition and displayed Cut case.

  $n = 0$ হলে $\pi$ একটি স্বতঃসিদ্ধ। সেটি
  $\lfalse \Sequent \quad$ স্বতঃসিদ্ধ হলে
  $\delta$ হলো $x:\lfalse \Sequent \lfalse$
  স্বতঃসিদ্ধ; অর্থাৎ $\Gamma' = \{x:\lfalse\}$
  এবং $!C \ident \lfalse$। $\pi$ যদি
  $!A \Sequent !A$ রূপের স্বতঃসিদ্ধ হয়,
  তাহলে $\delta$ হলো $x: !A \Sequent !A$।

  $n>0$ হলে $\pi$-র শেষে কোনো বিধি~$R$ আছে।
  আরোহের অনুমানে তার পূর্বধারণাগুলির
  অনুরূপ সিকোয়েন্টের !!{proof}s \Log{N2i}-তে
  আছে। দরকারে চিহ্ন বদলে ধরে নিতে পারি,
  ওই প্রমাণগুলিতে সব অবমুক্তি-চিহ্ন আলাদা;
  কোনো চিহ্ন~$x$ যদি অবমুক্ত হয়, তবে যে
  অনুমানবিধি তাকে অবমুক্ত করে তার উপরের
  অংশেই সেটি ঘটে। এবার দেখাই, ওই !!{proof}s
  মিলিয়ে উপযুক্ত $\Gamma' \Sequent !C$-এর
  প্রমাণ পাওয়া যায়। $R$ অনুযায়ী ক্ষেত্রগুলি:
  % BN-SRC-560 TARGET-CORRECTION-BEGIN
  \par\smallskip\noindent\textnormal{\small\textbf{উৎস-সংশোধন (BN-SRC-560):} এই আরোহবাক্যে উৎসে N2c লেখা, যদিও প্রস্তাবনা ও সব ক্ষেত্র N2i-র; এখানে N2i করা হয়েছে।}
  % BN-SRC-560 TARGET-CORRECTION-END

  \begin{enumerate}
    \item $R$ যদি \RightR{\Weakening} হয়, $\pi$-র শেষাংশ
    \[
    \AxiomC{}
    \RightLabel{$\pi_1$}
    \Deduce$\Gamma \fCenter$
    \RightLabel{\RightR{\Weakening}}
    \UnaryInf$\Gamma \fCenter !A$
    \DisplayProof
    \]
    $\pi_1$-এ আরোহের অনুমান প্রয়োগে
    $\Gamma' \Sequent \lfalse$-এর !!a{proof}
    পাই। তাতে $\FalseInt$ প্রয়োগে
    $\Gamma' \Sequent !A$-এর !!a{proof} পাই।
    \item $R$ যদি \LeftR{\Weakening} হয়,
    $\pi$-র শেষাংশ
    \[
    \AxiomC{}
    \RightLabel{$\pi_1$}
    \Deduce$\Gamma \fCenter \Delta$
    \RightLabel{\LeftR{\Weakening}}
    \UnaryInf$!A, \Gamma \fCenter \Delta$
    \DisplayProof
    \]
    $\pi_1$-এ আরোহের অনুমান প্রয়োগে
    $\Gamma' \Sequent !C$-এর !!a{proof}
    পাই; সেটিই $\delta$ নিই। কারণ
    $\Gamma'$ যদি $\Gamma$-র সঙ্গত হয়,
    তবে $!A,\Gamma$-রও সঙ্গত: $\Gamma'$-তে
    $x : !A$-এর সংখ্যা $\Gamma$-তে~$!A$-এর
    অনুলিপির সংখ্যার বেশি নয়, তাই
    $\Gamma \cup \{!A\}$-তে~$!A$-এর
    অনুলিপির সংখ্যারও বেশি নয়।
    % BN-SRC-561 TARGET-CORRECTION-BEGIN
    \par\smallskip\noindent\textnormal{\small\textbf{উৎস-সংশোধন (BN-SRC-561):} উৎসচিত্রে বাঁ-দুর্বলীকরণের বিধি ডান-দুর্বলীকরণ নামে চিহ্নিত; প্রদর্শিত বাম অংশে !A যোগের সঙ্গে মিলিয়ে \LeftR{\Weakening} করা হয়েছে।}
    % BN-SRC-561 TARGET-CORRECTION-END
    \item $R$ যদি \LeftR{\Contraction} হয়,
    $\pi$-র শেষাংশ
    \[
    \AxiomC{}
    \RightLabel{$\pi_1$}
    \Deduce$!A, !A, \Gamma \fCenter \Delta$
    \RightLabel{\LeftR{\Contraction}}
    \UnaryInf$!A, \Gamma \fCenter \Delta$
    \DisplayProof
    \]
    $\pi_1$-এ আরোহের অনুমান প্রয়োগে
    $\Pi, \Gamma' \Sequent !C$-এর
    !!a{proof} পাই, যেখানে
    $\Pi \subseteq \{x : !A, y:!A\}$।
    $\Pi = \{x : !A, y:!A\}$ হলে,
    ওই N2 উপপ্রমাণ~$\delta_1$-এ
    $y:!A$ রূপের সব চিহ্নযুক্ত
    !!{formula}s-কে $x:!A$ দিয়ে বদলাই।
    $x$ উপপ্রমাণ~$\delta_1$-এর অন্তিম
    সিকোয়েন্টের প্রসঙ্গে থাকে বলে
    ধরে নিতে পারি, $\delta_1$-এর কোনো
    অনুমানবিধি $x:!B$ রূপের !!a{formula}
    অবমুক্ত করে না; অর্থাৎ $\delta_1$-এর
    কোনো সিকোয়েন্টের বাঁ দিকে $x:!B$
    সক্রিয় নয়। অতএব ফলটিও \Log{N2i}-তে
    !!a{proof}।
    % BN-SRC-562 TARGET-CORRECTION-BEGIN
    \par\smallskip\noindent\textnormal{\small\textbf{উৎস-সংশোধন (BN-SRC-562):} সংকোচনের চিহ্ন-বদল ও অবমুক্তি-পরীক্ষা উৎসে G2 উপপ্রমাণ $\pi_1$-তে করা হয়েছে; চিহ্নযুক্ত সিকোয়েন্টগুলি আসলে আরোহফলে পাওয়া N2 উপপ্রমাণ $\delta_1$-তে, তাই তিন স্থানে সেটিই নেওয়া।}
    % BN-SRC-562 TARGET-CORRECTION-END
    \item $R$ যদি $\RightR{\lif}$ হয়,
    $\pi$-র শেষাংশ
    \[
    \AxiomC{}
    \RightLabel{$\pi_1$}
    \Deduce$!A, \Gamma \fCenter !B$
    \RightLabel{\RightR{\lif}}
    \UnaryInf$\Gamma \fCenter !A \lif !B$
    \DisplayProof
    \]
    আরোহের অনুমানে $\Gamma$-র সঙ্গত
    $\Gamma'$-র জন্য $x: !A, \Gamma'
    \Sequent !B$ অথবা $\Gamma' \Sequent !B$-এর
    !!a{proof}~$\delta_1$ আছে। $\delta_1$-এর
    পরে \Intro{\lif} প্রয়োগ করে
    $\delta$ পাই।
    % BN-SRC-563 TARGET-CORRECTION-BEGIN
    \par\smallskip\noindent\textnormal{\small\textbf{উৎস-সংশোধন (BN-SRC-563):} উৎসের শেষ বাক্যে নির্মাণের দিক উল্টো: $\delta_1$-কে $\delta$ থেকে নয়, $\delta_1$-এর পরে →-ভূমিকা দিয়ে $\delta$ গড়া হয়।}
    % BN-SRC-563 TARGET-CORRECTION-END
    \item $R$ যদি \LeftR{\lif} হয়,
    $\pi$-র শেষাংশ
    \[
    \AxiomC{}
    \RightLabel{$\pi_1$}
    \Deduce$\Gamma_1 \fCenter !A$
    \AxiomC{}
    \RightLabel{$\pi_2$}
    \Deduce$!B, \Gamma_2 \fCenter \Delta$
    \RightLabel{\LeftR{\lif}}
    \BinaryInf$!A \lif !B, \Gamma_1, \Gamma_2 \fCenter \Delta$
    \DisplayProof
    \]
    আরোহের অনুমানে $\Gamma_1' \Sequent !A$-এর
    !!{proof}s~$\delta_1$ এবং
    $x: !B, \Gamma_2' \Sequent !C$ অথবা
    $\Gamma_2' \Sequent !C$-এর
    $\delta_2$ পাই। $\delta_2$ দ্বিতীয় রূপের
    হলে $\delta_2$-ই $\delta$ হতে পারে। নইলে
    N2 উপপ্রমাণ~$\delta_1$-এর সব সূত্রচিহ্ন
    ও অবমুক্তি-চিহ্ন দরকারে বদলে নিই,
    যাতে $\delta_1$ ও~$\delta_2$-তে
    কোনো অভিন্ন চিহ্ন না থাকে। তখন
    $\Gamma_1' \cup \Gamma_2'$ হলো
    $\Gamma_1 \cup \Gamma_2$-র সঙ্গত।
    $x:!B$ যেহেতু~$\delta_2$-এর অন্তিম
    সিকোয়েন্টে আছে, তাই~$\delta_2$-তে
    $x$ অবমুক্তি-চিহ্নের কোনো অনুমানবিধিতে
    $x:!B$ সক্রিয় হতে পারে না। $\delta_2$-র
    প্রতিটি $x:!B \Sequent !B$ স্বতঃসিদ্ধের
    স্থলে !!{proof}~$\delta_3$ বসাই:
    \[
    \Axiom$x: !A \lif !B \fCenter !A \lif !B$
    \AxiomC{}
    \RightLabel{$\delta_1$}
    \Deduce$\Gamma_1' \fCenter !A$
    \RightLabel{\Elim{\lif}}
    \BinaryInf$x: !A \lif !B, \Gamma_1' \fCenter !B$
    \DisplayProof
    \]
    আর $\delta_2$-তে $x:!B$-এর প্রতিটি
    সংঘটন $x:!A \lif !B$ দিয়ে বদলাই এবং প্রতিস্থাপিত
    স্বতঃসিদ্ধের নিচের প্রসঙ্গগুলিতে $\Gamma_1'$ যোগ করি।
    ফল $\Subst{\delta_2}{\delta_3}{x:!B}$
    হলো $x:!A \lif !B, \Gamma_1',
    \Gamma_2' \Sequent !C$-এর !!a{proof}।
    % BN-SRC-564 TARGET-CORRECTION-BEGIN
    \par\smallskip\noindent\textnormal{\small\textbf{উৎস-সংশোধন (BN-SRC-564):} বাঁ → ক্ষেত্রের চিহ্ন-বদল উৎসে G2 উপপ্রমাণ $\pi_1$-তে নির্দেশিত; N2 চিহ্নযুক্ত উপপ্রমাণ $\delta_1$-তেই তা প্রযোজ্য।}
    % BN-SRC-564 TARGET-CORRECTION-END
    % BN-SRC-565 TARGET-CORRECTION-BEGIN
    \par\smallskip\noindent\textnormal{\small\textbf{উৎস-সংশোধন (BN-SRC-565):} উৎসের $\delta_3$ বৃক্ষের সিদ্ধান্তে $\Gamma_1'$ বাদ পড়ে, যদিও দ্বিতীয় শাখা $\Gamma_1' \Sequent !A$-এর উপর নির্ভর করে; N2 →-নিষ্কাশনে ওই প্রসঙ্গও সিদ্ধান্তে রাখা হয়েছে।}
    % BN-SRC-565 TARGET-CORRECTION-END
% BN-SRC-878 TARGET-CORRECTION-BEGIN
\par\smallskip\noindent\textnormal{\small\textbf{পাঠ-সংশোধন (BN-SRC-878):} কলম-স্থাপিত প্রমাণের দ্বিতীয় শাখার খোলা অনুমানও
প্রতিস্থাপিত স্বতঃসিদ্ধের নিচে রাখতে হবে। কেবল পুরোনো চিহ্নযুক্ত
সূত্রের নাম বদলালে ওই শাখার প্রসঙ্গ হারিয়ে যায়।}
% BN-SRC-878 TARGET-CORRECTION-END
    \item $R$ যদি \RightR{\land} হয়,
    $\pi$-র শেষাংশ
    \[
    \AxiomC{}
    \RightLabel{$\pi_1$}
    \Deduce$\Gamma_1 \fCenter !A$
    \AxiomC{}
    \RightLabel{$\pi_2$}
    \Deduce$\Gamma_2 \fCenter !B$
    \RightLabel{\RightR{\land}}
    \BinaryInf$\Gamma_1, \Gamma_2 \fCenter !A \land !B$
    \DisplayProof
    \]
    আরোহের অনুমানে $\Gamma_1' \Sequent !A$-এর
    প্রমাণ~$\delta_1$ এবং
    $\Gamma_2' \Sequent !B$-এর
    $\delta_2$ পাই। N2 উপপ্রমাণ~$\delta_2$-র
    সব সূত্রচিহ্ন ও অবমুক্তি-চিহ্ন বদলে
    $\delta_1$ ও~$\delta_2$-র চিহ্নগুলি
    পৃথক রাখা যায়। তখন
    $\Gamma_1' \cup \Gamma_2'$ হলো
    $\Gamma_1 \cup \Gamma_2$-র সঙ্গত।
    $\delta_1$ ও~$\delta_2$-র অন্তিম
    সিকোয়েন্টে $\Intro{\land}$ প্রয়োগে
    !!a{proof}~$\delta$ পাই।
    % BN-SRC-566 TARGET-CORRECTION-BEGIN
    \par\smallskip\noindent\textnormal{\small\textbf{উৎস-সংশোধন (BN-SRC-566):} ডান $\land$ ক্ষেত্রের দ্বিতীয় আরোহফল উৎসে ভুলে $\Gamma_2' \Sequent !A$; প্রদর্শিত দ্বিতীয় পূর্বধারণা অনুযায়ী $!B$। চিহ্ন-বদলও G2 উপপ্রমাণ $\pi_2$ নয়, N2 উপপ্রমাণ $\delta_2$-তে করা হয়েছে।}
    % BN-SRC-566 TARGET-CORRECTION-END
    \item $R$ যদি $\LeftR{\land}$ হয়,
    যেমন, $\pi$-র শেষাংশ
    \[
    \AxiomC{}
    \RightLabel{$\pi_1$}
    \Deduce$!A, \Gamma \fCenter \Delta$
    \RightLabel{\LeftR{\land}}
    \UnaryInf$!A \land !B, \Gamma \fCenter \Delta$
    \DisplayProof
    \]
    আরোহের অনুমানে $\Gamma$-র সঙ্গত
    $\Gamma'$-র জন্য $x: !A, \Gamma'
    \Sequent !C$ অথবা $\Gamma' \Sequent !C$-এর
    !!a{proof}~$\delta_1$ আছে। দ্বিতীয়
    রূপে $\delta$ হিসেবেই $\delta_1$ নিই।
    প্রথম রূপে লক্ষ করি, $x:!A$ অন্তিম
    সিকোয়েন্টে থাকায় $\delta_1$-এর
    $x$ অবমুক্তি-চিহ্নের কোনো বিধিতে
    $x:!A$ সক্রিয় নয়। প্রতিটি
    $x:!A \Sequent !A$ স্বতঃসিদ্ধের
    স্থলে !!{proof}~$\delta_2$ বসাই:
    \[
    \Axiom$x: !A \land !B \fCenter !A \land !B$
    \RightLabel{\Elim{\land}}
    \UnaryInf$x: !A \land !B \fCenter !A$
    \DisplayProof
    \]
    এবং $x:!A$-এর প্রতিটি সংঘটন
    $x:!A \land !B$ দিয়ে বদলাই; অর্থাৎ
    $\Subst{\delta_1}{\delta_2}{x:!A}$ নিই।
    এটি $x:!A \land !B, \Gamma'
    \Sequent !C$-এর !!a{proof}।
    % BN-SRC-567 TARGET-CORRECTION-BEGIN
    \par\smallskip\noindent\textnormal{\small\textbf{উৎস-সংশোধন (BN-SRC-567):} দ্বিতীয় রূপে উৎসে $\delta_1$-কে $\delta$ থেকে পাওয়া বলা হয়েছে; এখানে ইতিমধ্যে পাওয়া $\delta_1$-কেই চাওয়া $\delta$ নেওয়া হয়েছে।}
    % BN-SRC-567 TARGET-CORRECTION-END
    \item $R$ যদি \Cut হয়, $\pi$-র শেষাংশ
    \[
    \AxiomC{}
    \RightLabel{$\pi_1$}
    \Deduce$\Gamma_1 \fCenter !A$
    \AxiomC{}
    \RightLabel{$\pi_2$}
    \Deduce$!A, \Gamma_2 \fCenter \Delta$
    \RightLabel{\Cut}
    \BinaryInf$\Gamma_1, \Gamma_2 \fCenter \Delta$
    \DisplayProof
    \]
    আরোহের অনুমানে $\Gamma_1' \Sequent !A$-এর
    প্রমাণ~$\delta_1$ এবং
    $x:!A, \Gamma_2' \Sequent !C$ অথবা
    $\Gamma_2' \Sequent !C$-এর
    $\delta_2$ পাই। দ্বিতীয় রূপে
    $\delta$ হিসেবে~$\delta_2$ নিই।
    অন্যথায় $\delta_2$-র প্রতিটি
    $x:!A \Sequent !A$ স্বতঃসিদ্ধের
    স্থলে $\delta_1$ বসিয়ে
    $\Gamma_1', \Gamma_2' \Sequent !C$-এর
    !!{proof}~$\Subst{\delta_2}{\delta_1}{x:!A}$
    অর্থাৎ~$\delta$ পাই।
  \end{enumerate}
\end{proof}

\end{document}
